% !TeX program = xelatex
% CurveLab - standalone vector export. Compile with XeLaTeX, not pdfLaTeX.
% Structure: tailieuvehinh.pdf, I.1 (p.4); coordinates, draw, node, scope, clip.
% 3D is the current parallel projection, with visible/hidden line parts (I.3.5, pp.30-32).
% Captures the current viewport and sampled paths, without embedding a PNG.
% Width 16 cm; line widths and label sizes preserve the current drawing proportions.
% Mode: Hinh hoc phang
\documentclass[varwidth=50cm,border=2pt]{standalone}
\usepackage{fontspec}
\IfFontExistsTF{Times New Roman}{\setmainfont{Times New Roman}}{\setmainfont{TeX Gyre Termes}}
\IfFontExistsTF{Arial}{\setsansfont{Arial}}{\setsansfont{TeX Gyre Heros}}
\usepackage{amsmath,amssymb}
\usepackage{tikz,graphicx}
\usetikzlibrary{calc,angles,arrows.meta}
\definecolor{cl0}{RGB}{193,90,16}
\definecolor{cl1}{RGB}{120,59,176}
\definecolor{cl2}{RGB}{176,53,135}
\definecolor{cl3}{RGB}{8,126,98}
\definecolor{cl4}{RGB}{20,99,214}
% Change only this width; geometry, strokes and labels scale together.
\newcommand{\FigureWidth}{16cm}
\begin{document}
\begin{center}
\resizebox{\FigureWidth}{!}{%
\begin{tikzpicture}[x=1cm,y=-1cm,
 declare function={figwidth=16; figheight=9.333333;},
  s1/.style={draw=cl0,draw opacity=1,line width=0.853583pt,line cap=butt,line join=miter},
  s2/.style={draw=cl1,draw opacity=1,line width=0.853583pt,line cap=butt,line join=miter},
  s3/.style={draw=cl2,draw opacity=1,line width=0.853583pt,line cap=butt,line join=miter},
  s4/.style={draw=cl3,draw opacity=1,line width=0.474213pt,line cap=butt,line join=miter},
  s5/.style={draw=cl3,draw opacity=1,line width=0.853583pt,line cap=butt,line join=miter},
  s6/.style={anchor=base west,inner sep=0pt,outer sep=0pt,text=cl3,text opacity=1,font={\sffamily\fontsize{6.164764}{7.397717}\selectfont}},
  s7/.style={fill=cl4,fill opacity=1},
  s8/.style={draw=cl4,draw opacity=1,line width=0.616476pt,line cap=butt,line join=miter},
  s9/.style={anchor=base west,inner sep=0pt,outer sep=0pt,text=cl4,text opacity=1,font={\rmfamily\bfseries\fontsize{7.113189}{8.535827}\selectfont}},
  s10/.style={fill=cl2,fill opacity=1},
  s11/.style={draw=cl2,draw opacity=1,line width=0.616476pt,line cap=butt,line join=miter},
  s12/.style={anchor=base west,inner sep=0pt,outer sep=0pt,text=cl2,text opacity=1,font={\rmfamily\bfseries\fontsize{7.113189}{8.535827}\selectfont}},
  s13/.style={fill=cl3,fill opacity=1},
  s14/.style={draw=cl3,draw opacity=1,line width=0.616476pt,line cap=butt,line join=miter},
  s15/.style={anchor=base west,inner sep=0pt,outer sep=0pt,text=cl3,text opacity=1,font={\rmfamily\bfseries\fontsize{7.113189}{8.535827}\selectfont}}
]
% Coordinates use the displayed projection; original coordinates are retained in comments.
\path[use as bounding box] (0,0) rectangle ({figwidth},{figheight});
\begin{scope}
\clip (0,0) rectangle ({figwidth},{figheight});
% A; world = [0,0,0]
\coordinate (P1) at (5.171717,6.787879);
% B; world = [4,0,0]
\coordinate (P2) at (10.828283,6.787879);
% C; world = [0,3,0]
\coordinate (P3) at (5.171717,2.545455);
% I; world = [1,1,0]
\coordinate (P4) at (6.585859,5.373737);
% Drawing in display order: axes/grid, shapes, labels, formula.
\draw[s1] (P1)--(P2)--(P3)--(P1) (5.405051,6.787879)--(5.405051,6.554545)--(5.171717,6.554545);
\draw[s2] (2.626263,9.333333)--(11.959596,0) (5.605051,6.787879)--(5.604818,6.773701)--(5.604123,6.759537)--(5.602964,6.745405)--(5.601343,6.731317)--(5.599263,6.717291)--(5.596724,6.70334)--(5.593731,6.689479)--(5.590285,6.675724)--(5.586391,6.662089)--(5.582054,6.648588)--(5.577276,6.635237)--(5.572065,6.622049)--(5.566425,6.609039)--(5.560362,6.59622)--(5.553883,6.583607)--(5.546995,6.571212)--(5.539705,6.559049)--(5.532021,6.547132)--(5.523951,6.535472)--(5.515504,6.524082)--(5.506688,6.512975)--(5.497514,6.502162)--(5.487992,6.491655)--(5.47813,6.481466)--(5.467941,6.471604)--(5.457434,6.462082)--(5.446621,6.452908)--(5.435514,6.444092)--(5.424124,6.435645)--(5.412464,6.427575)--(5.400547,6.419891)--(5.388384,6.412601)--(5.375989,6.405713)--(5.363376,6.399234)--(5.350557,6.393171)--(5.337547,6.387531)--(5.324359,6.38232)--(5.311008,6.377542)--(5.297507,6.373205)--(5.283872,6.369311)--(5.270117,6.365865)--(5.256256,6.362872)--(5.242305,6.360333)--(5.228279,6.358253)--(5.214191,6.356632)--(5.200059,6.355473)--(5.185895,6.354777)--(5.171717,6.354545);
\draw[s3] (P4) ellipse [x radius=1.414141cm,y radius=1.414141cm];
\draw[s4] (P2)--(P1)--(P3);
\draw[s5] (5.405051,6.787879)--(5.405051,6.554545)--(5.171717,6.554545);
\node[s6] at (5.605051,7.187879) {90°};
\fill[s7] (P1) ellipse [x radius=0.083333cm,y radius=0.083333cm];
\draw[s8] (P1) ellipse [x radius=0.083333cm,y radius=0.083333cm];
\node[s9] at (5.321717,6.637879) {A};
\fill[s10] (P2) ellipse [x radius=0.083333cm,y radius=0.083333cm];
\draw[s11] (P2) ellipse [x radius=0.083333cm,y radius=0.083333cm];
\node[s12] at (10.978283,6.637879) {B};
\fill[s13] (P3) ellipse [x radius=0.083333cm,y radius=0.083333cm];
\draw[s14] (P3) ellipse [x radius=0.083333cm,y radius=0.083333cm];
\node[s15] at (5.321717,2.395455) {C};
\fill[s7] (P4) ellipse [x radius=0.083333cm,y radius=0.083333cm];
\draw[s8] (P4) ellipse [x radius=0.083333cm,y radius=0.083333cm];
\node[s9] at (6.735859,5.223737) {I};
\end{scope}
\end{tikzpicture}%
}
\end{center}
\end{document}
